\section{Conclusions and Future Work}
\label{sec:conclusion}

This paper presented four retrieval systems developed for the Scientific Retrieval task proposed by CLEF, addressing the challenge of matching 
informal social media queries to multilingual scientific documents in English, French, and German.
\noindent Our experiments expose a consistent progression in retrieval achievement (measured as scores of the metrics used) across the runs. The baseline, if on one hand established the upper bound for English queries, on the other hand highlighted the structural limits of an exact matching system based; indeed, the vocabulary mismatch between informal social media posts and scientific documents caused the collapse of the metrics scores for German and French.  The attempt to address this gap through the query expansion technique Pseudo Relevance Feedback has further decreased the scores, confirming the poor quality of the initial retrieval set. The introduction of hybrid dense-sparse retrieval through FAISS, combined 
with a Cross-Encoder re-ranker, produced a first substantial increment in the retrieval of non English queries; indeed, the deep neural networks used for German and French allowed the retrieval of documents thanks to the semantic embeddings similarity. 
\noindent The final architecture refined the hybrid pipeline by using a more powerful long-context retrieval model, which eliminated the previous chunking need. The results achieved were the highest and most balanced across all three languages, confirming that the hybrid dense-sparse retrieval with a neural re-ranking represents the most effective architecture for multilingual scientific document retrieval in this setting.